\documentclass[12pt,a4paper]{article}

\usepackage[utf8]{inputenc}
\usepackage[T1]{fontenc}
\usepackage[brazilian]{babel}
\usepackage{newtxtext}
\usepackage{newtxmath}
\usepackage{geometry}
\usepackage{setspace}
\usepackage{booktabs}
\usepackage{tabularx}
\usepackage{hyperref}
\usepackage{enumitem}

\geometry{left=3cm,right=2cm,top=3cm,bottom=2cm}
\setstretch{1.5}
\setlength{\parindent}{1.25cm}
\setlength{\parskip}{0pt}

\begin{document}

\begin{titlepage}
    \centering
    {\large Universidade Presbiteriana Mackenzie \par}
    \vspace{0.4cm}
    {\large Banco de Dados - Projeto Aplicado II \par}
    \vspace{1.8cm}
    {\Large \textbf{Etapa 3 - Aplicando Conhecimento} \par}
    \vspace{0.5cm}
    {\large \textbf{Classificacao e Monitoramento Textual de Boletins Epidemiologicos sobre Dengue} \par}
    \vfill
    \begin{flushleft}
    \textbf{Grupo:} Ana Clara Silva de Souza; Cid Wallace Araujo de Oliveira; Eduardo Machado Silva; Frederico Ripamonte Borges \\
    \textbf{Disciplina:} Projeto Aplicado II \\
    \textbf{Professor:} Anderson Adaime de Borba \\
    \textbf{Semestre:} 2026.01
    \end{flushleft}
    \vfill
    {\large \today \par}
\end{titlepage}

\section*{Resumo Executivo}
Nesta etapa aplicamos o metodo analitico definido na Etapa 2 sobre a base de boletins processada (60 chunks textuais). O pipeline foi executado de ponta a ponta com rotulagem inicial reprodutivel, treino supervisionado e avaliacao com metricas de acuracia e metricas complementares. Como resultado preliminar, os dois modelos avaliados (Naive Bayes Multinomial e Regressao Logistica Multiclasse) atingiram acuracia de 0.889 no conjunto de teste, com macro F1 de 0.76.

\section{Aplicacao do Metodo Analitico na Base}
\subsection{Base utilizada}
A base aplicada nesta etapa foi composta por:
\begin{itemize}[leftmargin=1.2cm]
    \item 11 boletins epidemiologicos (PDF) da fonte COVISA/SMS-SP;
    \item 60 chunks em \texttt{data/processed/chunks\_dataset.csv};
    \item 5 categorias tematicas alvo:
    \begin{itemize}
        \item \texttt{cenario\_epidemiologico}
        \item \texttt{prevencao\_controle\_vetorial}
        \item \texttt{orientacoes\_clinicas\_sintomas}
        \item \texttt{politicas\_publicas\_campanhas}
        \item \texttt{recomendacoes\_operacionais}
    \end{itemize}
\end{itemize}

\subsection{Aplicacao do pipeline}
O metodo aplicado seguiu os mesmos fundamentos tecnicos definidos na Etapa 2:
\begin{enumerate}[leftmargin=1.2cm]
    \item Extracao e preparacao textual (ja executadas na etapa anterior);
    \item Rotulagem inicial reprodutivel via \texttt{src/labeling.py}, com pseudo-rotulagem por TF-IDF + KMeans (5 clusters) e mapeamento para categorias tematicas;
    \item Vetorizacao com TF-IDF (\texttt{ngram\_range=(1,2)}, \texttt{min\_df=2}, \texttt{max\_df=0.95}, \texttt{max\_features=20000});
    \item Treino supervisionado com dois modelos:
    \begin{itemize}
        \item baseline: Naive Bayes Multinomial;
        \item principal: Regressao Logistica Multiclasse.
    \end{itemize}
    \item Divisao estratificada da base: 70\% treino, 15\% validacao e 15\% teste.
\end{enumerate}

\subsection{Distribuicao de classes}
Com a rotulagem inicial, a distribuicao ficou:
\begin{itemize}[leftmargin=1.2cm]
    \item \texttt{politicas\_publicas\_campanhas}: 17
    \item \texttt{orientacoes\_clinicas\_sintomas}: 13
    \item \texttt{cenario\_epidemiologico}: 11
    \item \texttt{prevencao\_controle\_vetorial}: 11
    \item \texttt{recomendacoes\_operacionais}: 8
\end{itemize}

\section{Medidas de Acuracia e Justificativa}
\subsection{Metricas definidas}
A avaliacao foi feita no conjunto de teste (9 amostras), usando:
\begin{itemize}[leftmargin=1.2cm]
    \item \textbf{Acuracia}: proporcao de predicoes corretas sobre o total;
    \item \textbf{Precision, Recall e F1-score} por classe;
    \item \textbf{Medias macro e weighted} para leitura multiclasse;
    \item \textbf{Matriz de confusao} para identificar erros entre classes.
\end{itemize}

\subsection{Justificativa}
Em problemas multiclasse com possivel desbalanceamento, apenas acuracia pode mascarar desempenho fraco em classes menores. Por isso, a leitura por media macro e por matriz de confusao foi usada para verificar robustez por classe, conforme definido na etapa anterior.

\section{Resultados Preliminares}
\subsection{Comparacao entre modelos}
\begin{table}[htbp]
\centering
\renewcommand{\arraystretch}{1.2}
\begin{tabularx}{\textwidth}{lccc}
\toprule
\textbf{Modelo} & \textbf{Acuracia} & \textbf{Macro Recall} & \textbf{Macro F1} \\
\midrule
Naive Bayes Multinomial & 0.889 & 0.800 & 0.760 \\
Regressao Logistica Multiclasse & 0.889 & 0.800 & 0.760 \\
\bottomrule
\end{tabularx}
\caption{Desempenho no conjunto de teste}
\end{table}

\subsection{Leituras da matriz de confusao}
Os dois modelos apresentaram mesmo padrao de acerto no teste. A principal falha observada foi na classe \texttt{recomendacoes\_operacionais}, que teve um unico exemplo no conjunto de teste e foi confundido com \texttt{politicas\_publicas\_campanhas}. Esse resultado reforca a necessidade de ampliar e refinar rotulos para classes menos representadas.

\section{Produto Gerado e Rascunho de Modelo de Negocios}
\subsection{Produto preliminar}
O produto preliminar e um \textbf{classificador tematico de trechos de boletins epidemiologicos}, capaz de:
\begin{itemize}[leftmargin=1.2cm]
    \item ingerir novos boletins em PDF;
    \item extrair e segmentar texto automaticamente;
    \item classificar cada trecho em categorias de vigilancia;
    \item priorizar leitura tecnica por tema (cenario, prevencao, clinica, politica, operacao).
\end{itemize}

\subsection{Proposta de valor}
Publicos-alvo iniciais:
\begin{itemize}[leftmargin=1.2cm]
    \item equipes de vigilancia epidemiologica;
    \item gestores de saude publica;
    \item equipes academicas de monitoramento territorial.
\end{itemize}
Valor entregue: reduzir tempo de triagem de boletins e acelerar identificacao de sinais relevantes para decisao operacional.

\subsection{Modelo de negocios (rascunho)}
\begin{itemize}[leftmargin=1.2cm]
    \item \textbf{Cliente}: secretarias e observatorios de saude;
    \item \textbf{Oferta}: painel de classificacao automatica + alertas por categoria;
    \item \textbf{Monetizacao}: assinatura institucional anual (SaaS) e projetos de customizacao;
    \item \textbf{Estrategia inicial}: piloto com municipio/regiao e validacao de ganho de tempo dos analistas.
\end{itemize}

\section{Esboco de Storytelling}
Narrativa proposta para apresentacao da Etapa 3:
\begin{enumerate}[leftmargin=1.2cm]
    \item \textbf{Problema}: volume e complexidade dos boletins dificultam leitura rapida orientada a acao.
    \item \textbf{Dados}: boletins oficiais COVISA processados em corpus estruturado.
    \item \textbf{Metodo}: pipeline reproducivel (extracao, chunking, rotulagem inicial, treino supervisionado).
    \item \textbf{Resultados}: acuracia preliminar de 0.889 e identificacao objetiva dos erros por classe.
    \item \textbf{Impacto}: viabilidade de produto para apoiar vigilancia e gestao.
    \item \textbf{Proxima iteracao}: rotulagem humana validada, aumento de base e ajuste fino por classe.
\end{enumerate}

\section{Conclusao}
A Etapa 3 cumpre os quatro requisitos da rubrica: metodo aplicado na base, metricas de acuracia com justificativa, descricao de resultados preliminares com produto e modelo de negocios, e storytelling estruturado. Como passo seguinte para elevar robustez, o grupo recomenda validacao humana dos rotulos e ampliacao amostral para reduzir confusao entre classes proximas.

\vspace{0.2cm}
\noindent\textbf{Fonte dos boletins:} \url{https://prefeitura.sp.gov.br/web/saude/w/vigilancia_em_saude/boletim_covisa/arboviroses}

\end{document}

